\section{City Bike Dataset}
\begin{comment}
    Description of the dataset, including its source
    Link to the notebook where the dataset was explored and analyzed (src/notebooks/...)  
\end{comment}

\subsection{Dataset Structure}
\begin{comment}
    What data is included, how it is organized (various csv files), variables, etc.
    Take in consideration also the features that we extracted from the original dataset 
    (trip duration, distance, day of the week, etc. that can be found in src/backend/services/historical.py)
\end{comment}
\subsubsection{Original Structure}
\subsubsection{Extracted Features}
\begin{comment}
    Take in consideration the choices made for computing the distance (Haversine distance vs Manhattan distance)
    1) Haversine distance is more accurate for long distances gives a straight line distance between two points, correct for the curvature of the Earth
       but might be computationally expensive and underestimate the actual distance traveled along the streets
    2) Manhattan distance is simpler and more appropriate for a city like New York with a grid-like street layout, it gives the distance traveled along the streets rather than a straight line distance
\end{comment}

\subsection{Analysis and Insights}
\begin{comment}
    Extracted insights from the dataset, also using visual features
\end{comment}
\subsection{Data Quality}
\begin{comment}
    Discussion of any data quality issues, such as missing values, outliers, or inconsistencies, and how they were addressed in the analysis.
    UUID wrongly formatted in the downloaded dataset, and describe the approach we applied to fix it
\end{comment}


